% !TeX root = ../main.tex

\chapter{研究设计与路径定义}
\section{基础复现与同规则延展}
基础成熟期要求最近20季齐全、成长率全部为正，当前成长率居全截面有效成长率前30\%；估值在全截面有限PE中取最低30\%后要求为正，最后取交集。2003—2007年采用自首条成长记录起每期为正且无断档的初始替代要求，不预先限制成长率百分比，在最终交集内按成长率降序取最多15只；同值以股票代码升序决定先后。成熟期交集全部纳入，不设15只上限。

季度末名单从次月起使用三个月。因此，2007年12月初始期名单仍影响2008年第一季度收益，成熟期持仓从2008年4月开始。基础延展只延长计算终点，不改变上述规则。

\section{四条优化路径}
\begin{table}[htbp]
\centering
\caption{六个正式方案的规则关系}
\renewcommand{\arraystretch}{1.18}
\small
\setlength{\tabcolsep}{4.5pt}
\begin{tabularx}{\textwidth}{@{}p{2.5cm}*{4}{C}@{}}
\toprule
方案 & 成熟期历史门槛 & 成长率与PE排名 & 成熟期组合 & 动态仓位 \\
\midrule
B 基础复现 & 20/20为正 & 全截面分别排名 & 交集全部纳入 & 无 \\
E 基础延展 & 同B & 同B & 同B & 无 \\
P1 路径一 & 至少19/20为正 & 历史合格候选池内 & 两池合并打分，最多15只 & 无 \\
P2 路径二 & 同P1 & 全截面分别排名 & 同P1 & 无 \\
P3 路径三 & 同P1 & 同P1 & 名单同P1 & 原有风控 \\
P4 路径四 & 同P2 & 同P2 & 名单同P2 & 原有风控 \\
\bottomrule
\end{tabularx}
\end{table}

四路径均恢复正PE条件及基础版初始五年规则，初始名单不足15只时不按三因子补位。P1是正式的“成功选股法优化”，P2是第二条独立排名路径；P3、P4分别叠加在P1、P2上，不是在两种选股规则之间再次择优切换。

P1、P2成熟期的核心候选要求两项名次均达到30\%阈值，扩展候选使用50\%阈值。核心与补位合并后统一按三维得分截取最多15只，核心只是候选类别，没有优先保留权。因此，本文不写“原筛选核心全部保留”。E到P1同时包含历史门槛、排名、组合容量与三因子得分等变化，其差异只能作为整套组合规则的比较。

\begin{figure}[htbp]
\centering
\includegraphics[width=\linewidth]{fig01}
\caption{基础规则与四条优化路径的对照关系}
\par\smallskip
\begin{minipage}{0.97\linewidth}\small 注：箭头表示设计上的对照关系，不表示收益因果分解。资料来源：CSMAR，作者计算。\end{minipage}
\end{figure}

\section{排名母体的精确定义}
基础版分别用当季所有有限成长率股票建立成长率榜，用所有有限PE股票建立估值榜，记样本数为Ng、Np。先排序，再与历史条件取交集；两榜样本数不必相等。P2恢复该口径，核心阈值为$\lfloor 0.3N_g\rfloor$和$\lfloor 0.3N_p\rfloor$，扩展阈值相应为50\%。

P1则先筛出同时具备有限成长率、有限PE且历史合格的H只股票，再在其中建立两榜；两个30\%阈值都基于H。它的准确含义是“合格候选池内前30\%”，不能称为“全市场前30\%”。排名采用离散名次及向下取整，不是分位数插值。

以2024年12月末为例，Ng=3,974、Np=5,121、H=82，对应30\%截止名次分别为1,192、1,536和24。名次与分母同时变化，因此两种规则既可能互相纳入，也可能互相排除股票，不能预先认定全截面排名一定更严格或更宽松。P1与P2的历史合格池及三因子得分已逐条核对一致，以避免混入新的评分变化。

\section{排名分母如何改变实际入选}
以当前已核验的2024年末为例，全截面成长率榜有3,974只、PE榜有5,121只，历史合格且两指标有限的候选池仅82只。基础版及P2核心的两项截止名次分别为1,192和1,536；P1在82只内部的截止名次均为24。这里比较的是不同榜单上的位置，不能直接用“24小于1,192”推断某条规则更苛刻，因为单只股票的名次也随母体改变。

可以用假设例子理解：一只股票在合格池的成长率排第20、PE排第20，会通过P1两项核心阈值；若它在全市场成长率排第1,300，则不通过P2成长阈值。另一只股票在合格池成长率排第30，但在全市场排第1,000，且PE同时过关，则可能通过P2却不能通过P1。示例只说明逻辑可能性，实际组合变化仍以逐季名单核验为准。

负PE占位是另一个独立问题。若全市场有100条有限PE，其中25条为负，按基础规则先升序取最低30条，再要求PE为正，则正PE最多只剩5条，而不是正PE股票75只内部的最低30\%即22条。这不是说负PE股票可以买入，而是负PE参与构造比较门槛。把正PE筛选移动到排名之前，会改变阈值及名单，属于另一个实验；当前基础版与正式优化路径均不作这种隐含替换。

P1到P2的设计对照固定历史候选条件与得分，再改变成长和PE的排名母体，因而能更集中地观察排名实施效果。E到P1则还包含20/20变19/20、扩展池、三维得分和成熟期持仓上限变化；即使收益提高，也无法仅从这一组差额分离出某一个环节的贡献。

\section{动态仓位的实际机制}
风控首先用此前最多12个月底层未控净收益计算年化波动率，至少6个历史月后实施20\%目标波动率。仓位为目标波动除以历史波动并截于1，无杠杆；历史不足或波动不可用时维持满仓。

\begin{equation}
w_t=\min\left(1,\frac{0.20}{\hat{\sigma}_{t-1}}\right).
\label{eq:weight}
\end{equation}

另在每个新季度将参考净值重置为该季度起点净值。月初根据已完成的净值判断：若相对季度起点下跌达到25\%，该季度余下月份仓位置0；下一季度重置。代码未在季度内持续更新峰值，因此它是季度起点损失阈值，不能称为滚动最高净值回撤止损，也不能保证月内跌幅不超过25\%。

实际受控收益为底层已扣基础费用的净收益乘以仓位。现金收益记0，未另计动态仓位变化成本；基础费用也随仓位缩放，不是逐笔交易成本模型。P3和P4不改变持仓名单，减少的是风险暴露而非股票数量。

\section{收益实现与评价指标}
每月对有收益记录的目标持仓重新等权平均；当月没有有效收益则记0。每个有持仓季度首月固定扣0.13\%，不按真实换手率计算。该实现含月度再等权和缺失收益再平均，不是季度买入后权重自然漂移的严格模拟。

\begin{equation}
CAGR=\left[\prod_{t=1}^{T}(1+R_{p,t})\right]^{12/T}-1.
\label{eq:cagr}
\end{equation}

波动为月收益样本标准差乘$\sqrt{12}$；最大回撤包含初始净值1；夏普沿用程序零无风险利率口径；Calmar为CAGR除以最大回撤绝对值。所有跨方案主表统一月份。本文没有为多路径比较新增显著性检验，也不将较高夏普或CAGR差直接解释为独立alpha。历史19/20参数参考过后续表现，所有分段均为回溯描述，不称为未接触的样本外测试。
